\documentclass[../main.tex]{subfiles}

\begin{document}

\section{Chapter 2: Data Types, Marks, and Channels}
\label{sec:chapter2-marks-channels}

\subsection{Techniques / Principles Applied}

\begin{itemize}[leftmargin=*]
\item Classifying variables as nominal, ordinal, or quantitative before
encoding them.
\item Selecting marks that fit the analytical task: cells, bars, flows, and
lines.
\item Prioritizing precise channels such as aligned position and length for the
most important quantitative comparisons.
\item Using color, width, and enclosure as supporting channels rather than the
primary source of exact comparison.
\end{itemize}

\subsection{How Applied in the Dashboard}

\subsubsection{Classifying Variables Before Encoding}

The dashboard first separates variables by role before assigning any visual
form. Income quintile and education in Figure~\ref{fig:socioeconomic-matrix}
are treated as ordered categories, payment channels in
Figure~\ref{fig:national-inflows} are treated as nominal classes within each
inflow type, and rates or weighted shares across
Figures~\ref{fig:mobile-ownership-profile} and
\ref{fig:save-borrow-trajectories} are treated as quantitative values. That
classification step determines whether the page should emphasize comparison,
composition, structure, or directional change.

Table~\ref{tab:marks-channels-summary} summarizes the main encodings used
across the five pages.

\begin{table}[H]
\centering
\small
\begin{tabularx}{\textwidth}{p{0.16\textwidth}p{0.23\textwidth}p{0.24\textwidth}X}
\toprule
\textbf{Chart/Page} & \textbf{Variables} & \textbf{Marks / channels} & \textbf{Why appropriate} \\
\midrule
Page 1 socioeconomic matrix & Income quintile, education, survey weight & Grid cells; row and column position; sequential color intensity & Two ordered categorical variables form a discrete matrix, while weighted population density is conveyed by ordered color. \\
Page 2 mobile ownership profile & Location, gender, income group, phone type, access rate & Bars; aligned length; grouped position; hue & Bar length supports accurate rate comparison across small multiples and hue separates phone types without changing the baseline. \\
Page 3 account composition flow & Location, account state, weighted volume & Sankey nodes and ribbons; horizontal position; ribbon width; color grouping & The chart emphasizes how groups move between account states, which is more important here than exact pairwise magnitude reading. \\
Page 4 inflow digitization & Inflow type, payment modality, percentage share & Segmented bars; aligned position; segment length; hue & A 100\% stacked bar makes composition comparable across wages, agriculture, transfers, and other inflows. \\
Page 5 save vs. borrow trajectories & Income group, income runway, save rate, borrow rate & Lines; vertical position; color & The line overlays help readers track directional change across ordered categories. \\
\bottomrule
\end{tabularx}
\caption{Main marks, channels, and encoding rationale used in the dashboard.}
\label{tab:marks-channels-summary}
\end{table}

\dashboardfigure{0.72}{fig-page1-socioeconomic-matrix.png}
{The socioeconomic matrix (income vs. education).}
{fig:socioeconomic-matrix}

\dashboardfigure{0.82}{fig-page2-mobile-ownership.png}
{Mobile ownership rate by location, phone type, gender, and income group.}
{fig:mobile-ownership-profile}

\dashboardfigure{0.72}{fig-page3-account-flow.png}
{The account composition flow.}
{fig:account-composition-flow}

\dashboardfigure{0.88}{fig-page4-inflows.png}
{The national digitization of financial inflows: cash vs. digital delivery.}
{fig:national-inflows}

\begin{figure}[H]
    \centering
    \begin{subfigure}[t]{0.48\textwidth}
        \centering
        \includegraphics[width=\textwidth]{fig-page5-save-borrow-income-group.png}
        \caption{By income group.}
        \label{fig:save-borrow-income-group}
    \end{subfigure}
    \hfill
    \begin{subfigure}[t]{0.48\textwidth}
        \centering
        \includegraphics[width=\textwidth]{fig-page5-save-borrow-income-runway.png}
        \caption{By income runway length.}
        \label{fig:save-borrow-income-runway}
    \end{subfigure}
\caption{Save vs. borrow percentage overlays.}
\label{fig:save-borrow-trajectories}
\end{figure}

\subsubsection{Selecting Marks for the Analytical Task}

The chosen marks follow the main question on each page. Figure~\ref{fig:socioeconomic-matrix}
uses matrix cells to compare ordered categorical intersections, Figure~\ref{fig:mobile-ownership-profile}
uses grouped bars for subgroup rate comparison, Figure~\ref{fig:account-composition-flow}
uses ribbons and nodes to emphasize movement between account states,
Figure~\ref{fig:national-inflows} uses segmented bars for within-category
composition, and Figure~\ref{fig:save-borrow-trajectories} uses line overlays
to highlight directional change across ordered groups.

\subsubsection{Prioritizing Position and Length for Exact Comparison}

Where precise reading matters most, the dashboard relies on aligned position
and length. Figure~\ref{fig:mobile-ownership-profile} makes subgroup gaps easy
to compare because all bars share common baselines, while
Figure~\ref{fig:national-inflows} lets readers compare the dominant payment
mode within each inflow category. Even in Figure~\ref{fig:socioeconomic-matrix},
ordered row and column position do the main structural work before color is
used to show density.

\subsubsection{Using Supporting Channels Carefully}

Supporting channels are kept secondary to the main quantitative comparison.
Figure~\ref{fig:socioeconomic-matrix} uses sequential blue intensity to signal
weighted density, Figure~\ref{fig:account-composition-flow} uses ribbon width
and grouped color to organize structural flows, and
Figure~\ref{fig:save-borrow-trajectories} uses hue to distinguish saving from
borrowing once the reader has already located the values by position.

\subsection{Notes / Adjustments}

Figure~\ref{fig:account-composition-flow} sacrifices some magnitude precision
because ribbon width is less exact than bar length, but it better communicates
movement between account states. Figure~\ref{fig:save-borrow-trajectories}
connects ordinal categories, so it implies continuity more than a strict
categorical display would; the choice is kept to make change in
save-versus-borrow behavior easier to track.

\end{document}
